% Generated draft: requires mathematical and source review.
\documentclass[11pt]{article}
\usepackage[margin=25mm]{geometry}
\usepackage{fontspec}
\setmainfont{Latin Modern Roman}
\usepackage{amsmath,amssymb,mathtools,unicode-math}
\setmathfont{Latin Modern Math}
\usepackage{enumitem,needspace,xurl,hyperref}
\setlength{\parindent}{0pt}
\setlength{\parskip}{5pt}
\setlength{\emergencystretch}{3em}
\newcommand{\N}{\mathbb{N}}
\newcommand{\Z}{\mathbb{Z}}
\newcommand{\Q}{\mathbb{Q}}
\newcommand{\R}{\mathbb{R}}
\newcommand{\C}{\mathbb{C}}
\newcommand{\F}{\mathbb{F}}
\newcommand{\A}{\mathbb{A}}
\newcommand{\PP}{\mathbb{P}}
\newcommand{\E}{\mathbb{E}}
\newcommand{\eps}{\varepsilon}
\newcommand{\dd}{\,\mathrm{d}}
\begin{document}
\section*{UnsolvedMath ID: 30005347}
\textbf{Source problem code:} catalogue entry
\subsection{Definitions and mathematical statement}
For a positive integer \(n\), let \(S_n\) be the group of permutations of \(\{1,\ldots,n\}\), let \(e\) denote its identity, and let
\[
T_n=\{(i\ j):1\leq i<j\leq n\}
\]
be its set of transpositions. A nonempty selection of pairwise distinct transpositions has the required property if its members admit an ordering \(\tau_1,\ldots,\tau_r\), with \(r\geq1\), such that
\[
\tau_1\tau_2\cdots\tau_r=e.
\]
The supplied target is to determine the least integer \(M_n\) such that every family of \(M_n\) transpositions in \(S_n\) contains such a selection.

\emph{Conditional precise formulation.} If ``family'' means a set of distinct transpositions, the intended selection must be nonempty, and the threshold ranges over positive integers, this becomes
\[
M_n=\min\Bigl\{m\in\Z_{\geq1}:\ \forall F\subseteq T_n,\ |F|=m\ \Longrightarrow\ \exists r\in\{1,\ldots,m\}\ \exists\tau_1,\ldots,\tau_r\in F,
\ \bigl(\tau_i\ne\tau_j\text{ for }i\ne j\bigr)
\ \land\ \tau_1\cdots\tau_r=e\Bigr\}.
\]
Determine \(M_n\) under the intended conventions.

\emph{Unresolved qualifications.} The record does not specify whether a family is a set or an indexed family allowing repetitions. With repetitions allowed but selected transpositions required to be distinct as group elements, arbitrarily large families consisting of one repeated transposition obstruct any finite threshold for \(n\geq2\). The set formulation above is therefore conditional, not an unqualified reconstruction. It also permits \(m>|T_n|=\binom n2\), where the universal condition is vacuous; the source does not specify whether such thresholds are admissible. Nonemptiness of the selection, positivity of the threshold, and the intended range of \(n\), including small boundary cases, require confirmation. Allowing an empty selection would make the identity-product condition automatic.

The supplied original excerpt refers to an unavailable preceding theorem and a question of Keevash, Mubayi, Sudakov, and Verstraete, with an unresolved reference label [1]. Its lower-bound typography is corrupted. The supplied curation interprets the relevant lower bound as \(\Omega(n\log n)\) and describes a near-matching upper bound from rainbow-cycle methods, with the exact value and a remaining logarithmic factor open. These are supplied contextual claims whose connection to the precise intended family convention requires review.
\subsection{Short English statement}
Determine the smallest family size forcing a nonempty collection of distinct transpositions that can be ordered to multiply to the identity. The meaning of ``family,'' the admissible thresholds, and the range of \(n\) need confirmation. The displayed formalization uses an explicitly conditional set interpretation. This is a statement-only draft requiring human mathematical review.
\subsection{Sources}
\begin{itemize}
\item Local registry: \texttt{data/unsolvedmath/problems.json}, permanent UnsolvedMath ID \(30005347\), title ``Additive Dimension of Transpositions,'' exact source code \texttt{catalogue entry}.
\item Supplied report citation: \emph{Combinatorics} (2023), Oberwolfach Reports, pp.\ 5--89, DOI \texttt{10.4171/owr/2023/1}. \url{https://doi.org/10.4171/owr/2023/1}. The supplied record describes Oberwolfach Reports as documenting mathematical workshops and their open-problem discussions.
\item Supplied literature source: ``Essentially Tight Bounds for Rainbow Cycles.'' \url{https://arxiv.org/abs/2309.04460}.
\item Supplied literature source: ``Essentially Tight Bounds for Rainbow Cycles, PLMS.'' \url{https://doi.org/10.1112/plms.70044}.
\item Attribution: the retained original passage associates the preceding theorem with a question of Keevash, Mubayi, Sudakov, and Verstraete, citing [1]. The bibliography entry for [1] and the preceding theorem are not supplied. No individual proposer is specified; the supplied proposed year is 2023.
\item Literature assessment dated exactly \texttt{2026-08-22}: the supplied status is ``partially solved.'' The supplied assessment states a lower bound \(\Omega(n\log n)\), a near-matching upper bound from rainbow-cycle methods, and an unresolved exact value and logarithmic factor. This is supplied dated curation derived from the OWR material, not independently verified here; the record requests independent verification of time-sensitive status and cited claims.
\item The registry labels the statement \texttt{corrected\_verified} and says it was checked against the official report. This records the supplier's verification claim; no source visit or independent verification is asserted in this draft.
\end{itemize}
\end{document}
